\section{Background and Related Work}
\subsection{Electrical Reconfiguration and Software-Defined Control}
Baran and Wu formulated distribution reconfiguration in terms of electrical
losses and load balance, establishing the relevance of switch state and radial
network structure \cite{baran1989}. The present work uses radiality only as an
operation-level topological constraint. It does not implement their optimization
procedure, solve an optimal power flow, or benchmark a new route-selection
method.

OpenFlow provides an influential example of a programmable interface between
network control and forwarding elements \cite{mckeown2008}. Applying this idea
to an energy system requires care over which plane is being programmed. In
conventional SDN, the data plane forwards communication traffic. In the project
terminology, the ``power plane'' instead denotes conducting electrical paths;
the interface described here is not an OpenFlow implementation for breakers.

Software-defined microgrid control is already an established research topic.
Wang et al. virtualize microgrid control functions as software services,
reducing dependence on dedicated controller hardware \cite{wang2020}. Li and
Du employ SDN and OpenFlow to support programmable communications for networked
microgrid formation, partition, and reconfiguration \cite{li2021}. These works
preclude a broad novelty claim that separating control from physical microgrid
operation is itself new. The narrower focus here is retaining an evolving
network-team command contract while producing an inspectable PSCAD switching
experiment and measurement record.

\subsection{Communications and Cyber-Physical Simulation}
Communication behavior can influence control outcomes, but a delayed timestamp
is not equivalent to a simulated communication network. Dorsch et al. compare
legacy and software-defined communications in distributed smart-grid control,
explicitly considering communication infrastructure and agent interaction
\cite{dorsch2017}. The long-delay case in this paper instead imposes a later
restoration time. It illustrates the electrical cost of waiting, without
measuring packet loss, routing convergence, or SDN latency.

Federated simulators address a different problem from batch automation. FNCS
couples power-system and communication simulators with synchronization
strategies for their exchanges \cite{ciraci2014}. HELICS supports coordinated
event-driven and time-series simulations, including iteration among federates
\cite{palmintier2017}. Neither framework is integrated into the current
implementation. Their relevance is methodological: a persistent application
connection does not supply the shared-time and causality mechanisms required
for closed-loop co-simulation.

More recent cyber-physical emulation also extends well beyond scheduled
switching. Haque et al. combine an operational-microgrid-derived power model
with a synthetic communication network, DNP3 exchanges, and structured threat
scenarios \cite{haque2025}. Their work provides a useful contrast to the
illustrative feeder and deterministic commands used here; this paper makes
no comparative cybersecurity or resilience-performance claim.

\begin{table}[tb]
\centering\small
\caption{Positioning against selected related work. Entries describe scope,
not a ranking or a feature-completeness assessment.}
\label{tab:related}
\begin{tabularx}{\textwidth}{@{}p{32mm}XX@{}}
\toprule
Work & Principal emphasis & Distinction of this study \\
\midrule
Baran--Wu \cite{baran1989} & Electrical reconfiguration objectives & Evaluates supplied routes; no optimizer \\
Wang et al. \cite{wang2020} & Virtualized microgrid controllers & Command-to-EMT experiment integration \\
Li--Du \cite{li2021} & SDN-enabled networked microgrids & JSON batch interface, not live OpenFlow \\
FNCS / HELICS \cite{ciraci2014,palmintier2017} & Synchronized simulator federation & Sequential automation of a single EMT run \\
Haque et al. \cite{haque2025} & Cyber-physical threat emulation & Illustrative switching and fault scenarios \\
\bottomrule
\end{tabularx}
\end{table}

\subsection{PSCAD as the Electrical Execution Layer}
PSCAD exposes Python automation for loading and running projects
\cite{pscadAutomation}. Its timed breaker component supports one or two
operations associated with a breaker control signal \cite{pscadTimedBreaker}.
The supplied fault models represent a fault using switched resistances rather
than a detailed arc characteristic \cite{pscadFaults}. These documented
facilities are sufficient for the deterministic switching cases examined here.
The project-specific contribution lies in translating commands, managing run
completion and result freshness, and connecting the results to research
artifacts; it is not a replacement for the underlying EMT solver.
